\section{项目概述}

本项目围绕“转座元件（TE）知识图谱与智能问答系统构建”展开，目标是从公开文献中抽取与转座元件相关的结构化知识，并将其组织为可检索、可视化和可问答的知识图谱系统。项目以 LINE-1 相关文献为切入点，通过文献抓取、信息抽取、图数据库建模、前端可视化和大模型问答五个环节，完成从原始摘要文本到交互式知识系统的完整流程。

系统面向的主要用户包括课程教师、同学以及希望快速浏览转座元件相关知识的初学者。相较于传统的关键词检索方式，本项目希望通过图谱化结构展示实体之间的关系，并结合自然语言问答接口，提高用户获取 TE 相关知识的效率。

\subsection{研究背景}

转座元件（Transposable Elements, TEs）是基因组中能够发生位置转移的 DNA 序列，在基因组进化、基因表达调控和多种疾病发生过程中具有重要作用。近年来，与 LINE-1、Alu、LTR 等转座元件相关的研究文献持续增长，但这些知识分散在大量论文摘要和研究报告中，普通的文献检索方式难以快速提炼出结构化关系。

知识图谱能够将“实体—关系—实体”的信息组织为图结构，使转座元件、疾病、生物学功能和文献之间的联系更加直观。若再结合大模型的自然语言理解和生成能力，便可以在本地知识库基础上提供更友好的问答入口。因此，本项目将 TE 文献挖掘、图数据库和 RAG 问答三者结合，形成一个面向课程作业的综合型生物信息学系统。

\subsection{系统总体流程}

本项目的整体流程可以概括为五个阶段。首先，从 PubMed 等公共文献数据库中抓取与 LINE-1 等转座元件相关的文献摘要及元数据。其次，利用规则方法和大模型辅助抽取摘要中的 TE、Disease、Function、Paper 等实体及其关系，并输出为结构化结果。随后，对抽取结果进行去重、别名统一和规范化处理，再导入 Neo4j 图数据库，构建知识图谱主存储层。

在数据层完成后，系统通过前端页面将知识图谱局部子图可视化展示出来，支持节点搜索、关系查看和中英文展示切换。最后，用户可以通过网页中的问答界面输入自然语言问题，后端结合 Neo4j 图数据库检索结果与大模型生成能力，返回带有参考文献的学术化答案。整个系统形成了“文献获取 $\rightarrow$ 信息抽取 $\rightarrow$ 图数据库构建 $\rightarrow$ 可视化展示 $\rightarrow$ 智能问答”的闭环。

\figplaceholder{系统总体流程图（建议后续替换为正式架构图）}{fig:overall}

